\documentclass[12pt]{article}

\usepackage[margin=1in]{geometry}
\usepackage[T1]{fontenc}
\usepackage{amsmath,amssymb,amsthm,mathtools}
\usepackage{newtxtext,newtxmath}
\usepackage{microtype}
\usepackage{lineno}
\usepackage{setspace}
\usepackage[hidelinks]{hyperref}

\setstretch{1.48}
\linenumbers
\setlength{\parindent}{1.2em}
\setlength{\emergencystretch}{2em}

\newtheorem{theorem}{Theorem}[section]
\newtheorem{lemma}[theorem]{Lemma}
\newtheorem{corollary}[theorem]{Corollary}
\newcommand{\C}{\mathbb C}
\newcommand{\R}{\mathbb R}
\newcommand{\ee}{\mathrm e}
\newcommand{\ii}{\mathrm i}

\title{\bfseries Supplementary Information for\\
One quantum return record permits every future return probability}
\author{Michael Haufschild}
\date{28 July 2026}

\begin{document}
\maketitle

\section{Binomial and parity-conditioned spectra}

Fix \(n\geq1\). For \(j=0,\ldots,n\), define
\[
 q_j=\frac{{n\choose j}}{2^n},\qquad
 r_j=q_j\bigl[1+(-1)^{n-j}\bigr].
\]

\begin{lemma}[Positive unit mass]
The lists \(q=(q_j)\) and \(r=(r_j)\) are nonnegative and each has total
mass one.
\end{lemma}

\begin{proof}
Nonnegativity of \(q\) is immediate. The factor multiplying \(q_j\) in
\(r_j\) is either zero or two. The binomial theorem gives
\[
\sum_{j=0}^{n}q_j=2^{-n}(1+1)^n=1.
\]
Moreover,
\[
\sum_{j=0}^{n}q_j(-1)^{n-j}
=2^{-n}(1-1)^n=0
\]
because \(n>0\). Hence \(\sum_jr_j=1\).
\end{proof}

For \(\theta\in[0,1]\), set
\[
 w_{\theta,j}=(1-\theta)q_j+\theta r_j.
\]

\begin{corollary}[Continuum of probability measures]
Every \(w_\theta\) is nonnegative and has total mass one.
\end{corollary}

\begin{proof}
This is convexity and the preceding lemma.
\end{proof}

\section{Shared nonnegative-Hamiltonian realization}

For \(T>0\), use the Hilbert space \(\C^{n+1}\), its standard basis
\(\{|j\rangle\}_{j=0}^{n}\), and the diagonal Hamiltonian
\[
 K_{n,T}|j\rangle=\frac{\pi j}{T}|j\rangle.
\]
For \(\theta\in[0,1]\), define
\[
 |\psi_\theta\rangle
 =\sum_{j=0}^{n}\sqrt{w_{\theta,j}}\,|j\rangle,
\quad
 |\psi_\theta(t)\rangle
 =\sum_{j=0}^{n}\sqrt{w_{\theta,j}}\,
   \ee^{\ii\pi jt/T}|j\rangle.
\]

\begin{theorem}[One Hamiltonian and normalized pure states]
\(K_{n,T}\) is self-adjoint and nonnegative. Every
\(|\psi_\theta\rangle\) and \(|\psi_\theta(t)\rangle\) has norm one, and
\[
\langle\psi_\theta|\psi_\theta(t)\rangle
=\sum_{j=0}^{n}w_{\theta,j}\ee^{\ii\pi jt/T}.
\]
\end{theorem}

\begin{proof}
The Hamiltonian is diagonal with real nonnegative entries. Both squared
norms equal \(\sum_jw_{\theta,j}=1\). Orthogonality of the standard basis
gives the displayed inner product.
\end{proof}

\section{Permutation-invariant qubit realization}

Let the computational basis of \(n\) qubits be indexed by subsets
\(S\subseteq\{1,\ldots,n\}\), with \(S\) recording excited sites. Define
\[
 H_{\rm u}|S\rangle=\frac{\pi |S|}{T}|S\rangle
 =\frac{\pi}{T}\sum_{a=1}^{n}
 |1\rangle\langle1|_a\,|S\rangle .
\]
For any permutation \(\sigma\) of the sites,
\(|\sigma(S)|=|S|\); hence \(H_{\rm u}\) is invariant under the full
permutation group.

For a probability list \(w=(w_0,\ldots,w_n)\), define
\[
 |\Psi_w\rangle
 =\sum_{S\subseteq\{1,\ldots,n\}}
 \sqrt{\frac{w_{|S|}}{{n\choose |S|}}}\,|S\rangle .
\]

\begin{theorem}[One uniform field]
The state \(|\Psi_w\rangle\) is normalized and, for every real \(t\),
\[
 \langle\Psi_w|\ee^{\ii tH_{\rm u}}|\Psi_w\rangle
 =\sum_{j=0}^{n}w_j\ee^{\ii\pi jt/T}.
\]
For every \(k\geq0\),
\[
 \langle\Psi_w|H_{\rm u}^{\,k}|\Psi_w\rangle
 =\sum_{j=0}^{n}w_j\left(\frac{\pi j}{T}\right)^k.
\]
\end{theorem}

\begin{proof}
There are \({n\choose j}\) subsets of cardinality \(j\). Grouping the
squared coefficients by cardinality gives
\[
 \|\Psi_w\|^2
 =\sum_{j=0}^{n}{n\choose j}
 \frac{w_j}{{n\choose j}}=\sum_jw_j=1.
\]
The Hamiltonian phase and its \(k\)th power depend only on \(|S|\).
Grouping both expectation values by the same cardinality shells cancels
the binomial multiplicity and gives the displayed formulas.
\end{proof}

\begin{corollary}[Homogeneous closed-interval ambiguity]
Taking \(w=w_{\sqrt p}\) gives, in one fixed permutation-invariant
noninteracting qubit Hamiltonian, normalized symmetric pure preparations
whose earlier return histories and prescribed finite energy moments are
common to arbitrary accuracy, while their return probability at \(T\)
ranges over every \(p\in[0,1]\).
\end{corollary}

\begin{proof}
Apply the theorem and the amplitude, moment and window identities proved
below. Site permutations preserve each equal-amplitude cardinality shell,
so every preparation is symmetric.
\end{proof}

\section{Fixed-protocol isospectral realization}

Fix the reference basis state \(|\phi\rangle=|0\rangle\in\C^{n+1}\). For every
\(\theta\in[0,1]\), there is a unitary \(U_\theta\) such that
\[
 U_\theta|\phi\rangle=|\psi_\theta\rangle.
\]
Indeed, each singleton is an orthonormal family; extend both to orthonormal
bases and map one basis to the other. Define
\[
 \widetilde K_{\theta}=U_\theta^\dagger K_{n,T}U_\theta,\qquad
 |\phi_\theta(t)\rangle=U_\theta^\dagger|\psi_\theta(t)\rangle.
\]

\begin{theorem}[One preparation, projector and spectrum]
Every \(\widetilde K_\theta\) is self-adjoint and nonnegative, and all are
unitarily equivalent to \(K_{n,T}\). The state
\(|\phi_\theta(t)\rangle\) has norm one and
\[
 \langle\phi|\phi_\theta(t)\rangle
 =\langle\psi_\theta|\psi_\theta(t)\rangle.
\]
For every integer \(k\geq0\),
\[
 \langle\phi|\widetilde K_\theta^{\,k}|\phi\rangle
 =\sum_{j=0}^{n}w_{\theta,j}\left(\frac{\pi j}{T}\right)^k.
\]
\end{theorem}

\begin{proof}
Taking adjoints shows
\[
 \widetilde K_\theta^\dagger
 =U_\theta^\dagger K_{n,T}^\dagger U_\theta
 =\widetilde K_\theta.
\]
Unitary conjugacy preserves nonnegativity and the spectrum. Norm and inner
product preservation give the survival identity. Induction gives
\[
 U_\theta\widetilde K_\theta^{\,k}
 =K_{n,T}^{\,k}U_\theta.
\]
Apply this to \(|\phi\rangle\), then evaluate the diagonal operator in
\(|\psi_\theta\rangle\).
\end{proof}

\section{Nearest-neighbour realization}

For \(0\leq\theta<1\), all \(w_{\theta,j}\) are strictly positive because
\[
w_{\theta,j}\geq(1-\theta)q_j>0.
\]
Define the Krylov family
\[
 |v_k\rangle=K_{n,T}^{\,k}|\psi_\theta\rangle,
 \qquad k=0,\ldots,n.
\]

\begin{lemma}[Cyclic Vandermonde basis]
The vectors \(|v_0\rangle,\ldots,|v_n\rangle\) are linearly independent.
\end{lemma}

\begin{proof}
If \(\sum_{k=0}^{n}c_k|v_k\rangle=0\), its \(j\)th coordinate is
\[
 \sqrt{w_{\theta,j}}\sum_{k=0}^{n}
 c_k\left(\frac{\pi j}{T}\right)^k=0.
\]
Strict positivity of \(w_{\theta,j}\) leaves a square Vandermonde system
at the distinct nodes \(\pi j/T\). Hence every \(c_k=0\).
\end{proof}

Apply Gram--Schmidt to this ordered basis and write the resulting
orthonormal basis as \(|b_0\rangle,\ldots,|b_n\rangle\). Because
\(\|v_0\|=1\), \(|b_0\rangle=|\psi_\theta\rangle\).

\begin{theorem}[Exact path Hamiltonian below the endpoint]
Let \(U_\theta|j\rangle=|b_j\rangle\) and
\(\widetilde K_\theta=U_\theta^\dagger K_{n,T}U_\theta\).
Then \(\widetilde K_\theta\) has the spectrum
\(\{0,\pi/T,\ldots,\pi n/T\}\), the fixed endpoint state is
\(|0\rangle\), and
\[
\langle i|\widetilde K_\theta|j\rangle=0
\qquad\text{if }|i-j|>1.
\]
Every adjacent path element is nonzero:
\[
 \langle a|\widetilde K_\theta|a+1\rangle\ne0,
 \qquad 0\leq a<n.
\]
Its fixed-endpoint survival amplitude is \(A_\theta(t)\).
\end{theorem}

\begin{proof}
Gram--Schmidt preserves every initial span, so
\[
\operatorname{span}\{b_0,\ldots,b_j\}
=\operatorname{span}\{v_0,\ldots,v_j\}.
\]
Since \(K_{n,T}v_k=v_{k+1}\),
\[
K_{n,T}\operatorname{span}\{b_0,\ldots,b_j\}
\subseteq\operatorname{span}\{b_0,\ldots,b_{j+1}\}.
\]
Consequently \(\langle b_i|K_{n,T}|b_j\rangle=0\) when
\(i>j+1\); self-adjointness gives the symmetric half. This is precisely
tridiagonality after conjugation by \(U_\theta\). Unitary conjugacy
preserves the spectrum, \(U_\theta|0\rangle=|\psi_\theta\rangle\), and the
fixed-endpoint survival amplitude.

For connectedness, write \(f_k=K_{n,T}^k|\psi_\theta\rangle\) and let
\(|b_k\rangle\) be the normalized Gram--Schmidt vectors. The leading
coefficient \(\langle b_k|f_k\rangle\) is nonzero; otherwise the
nonzero \(k\)th Gram--Schmidt residual would have zero norm. Expanding
\(f_a\) in the basis and using \(K_{n,T}f_a=f_{a+1}\), tridiagonality
eliminates all terms except
\[
 \langle b_{a+1}|f_{a+1}\rangle
 =\langle b_a|f_a\rangle
  \langle b_{a+1}|K_{n,T}|b_a\rangle .
\]
Both leading coefficients are nonzero, hence every adjacent matrix element
is nonzero. Self-adjointness gives the reverse orientation.
\end{proof}

Putting \(\theta=\sqrt p\) proves the nearest-neighbour realization for
every \(0\leq p<1\).

The endpoint \(\theta=1\) is noncyclic because half its weights vanish.
It nevertheless has a tridiagonal realization.

\begin{theorem}[Exact block-path Hamiltonian at perfect return]
For \(n=2q+1\), there is a unitary \(U_1\) with
\(U_1|0\rangle=|\psi_1\rangle\) such that
\[
 \widetilde K_1=U_1^\dagger K_{n,T}U_1
 =
 \begin{pmatrix}
   J_{\rm odd}&0\\
   0&D_{\rm even}
 \end{pmatrix}.
\]
Here \(J_{\rm odd}\) is tridiagonal on \(q+1\) sites and
\(D_{\rm even}\) is diagonal on \(q+1\) sites. Thus
\(\widetilde K_1\) is tridiagonal, has the complete spectrum
\(\{0,\pi/T,\ldots,\pi n/T\}\), and its fixed-endpoint survival
probability at \(T\) is one.
\end{theorem}

\begin{proof}
The nonzero parity weights occur exactly at the odd indices
\(1,3,\ldots,2q+1\), are positive and sum to one. Restrict the diagonal
Hamiltonian to those \(q+1\) levels. Their weighted Krylov orbit is a basis
by the same Vandermonde argument as above. Gram--Schmidt therefore produces
a tridiagonal \(J_{\rm odd}\) whose first vector is the normalized endpoint
state.

The complementary even-energy eigenvectors
\(|0\rangle,|2\rangle,\ldots,|2q\rangle\) are already an orthonormal
eigenbasis; in that order their restriction is the diagonal matrix
\(D_{\rm even}\). Concatenate the odd Krylov basis and the even eigenbasis.
The two spectral subspaces are orthogonal and invariant, so all cross
matrix elements vanish. The resulting block diagonal matrix is
tridiagonal, including at the interface. Basis concatenation is unitary,
preserves the full spectrum and maps the fixed first site to
\(|\psi_1\rangle\). Its survival probability is consequently the
parity-conditioned value, equal to one at \(T\).
\end{proof}

\begin{theorem}[Perfect return forces path disconnection]
Let \(n=2q+1\), let \(U|0\rangle=|\psi_1\rangle\), and put
\(\widetilde K=U^\dagger K_{n,T}U\). If \(\widetilde K\) is tridiagonal,
then
\[
 \langle a|\widetilde K|a+1\rangle=0
\]
for at least one \(a\in\{0,\ldots,n-1\}\).
\end{theorem}

\begin{proof}
First consider any tridiagonal operator \(L\) on a path basis
\(|0\rangle,\ldots,|n\rangle\) for which every forward matrix element
\(\langle a|L|a+1\rangle\) is nonzero. If \(Lv=\lambda v\) and
\(\langle0|v\rangle=0\), the zeroth row of the eigenvalue equation gives
\(\langle0|L|1\rangle\langle1|v\rangle=0\), hence
\(\langle1|v\rangle=0\). Inductively, after the coordinates through \(a\)
vanish, row \(a\), tridiagonality and the eigenvalue equation give
\[
 \langle a|L|a+1\rangle\langle a+1|v\rangle=0.
\]
Thus every coordinate of \(v\) vanishes. A connected tridiagonal path
therefore has no nonzero eigenvector orthogonal to its first endpoint.

For the stated \(\widetilde K\), take any even spectral basis vector
\(|2r\rangle\). Then \(v=U^\dagger|2r\rangle\) is a nonzero eigenvector of
\(\widetilde K\), while
\[
 \langle0|v\rangle
 =\langle\psi_1|2r\rangle=0
\]
because every even spectral weight of \(|\psi_1\rangle\) vanishes. Hence
\(\widetilde K\) cannot have all forward path couplings nonzero.
\end{proof}

\begin{corollary}[Closed-interval nearest-neighbour realization]
After choosing an odd \(n>m\) large enough for the requested finite-window
precision, one fixed endpoint state and projector, one common spectrum and
tridiagonal Hamiltonians realize every \(p\in[0,1]\), while matching all
moments through order \(m\) and the common earlier record.
The representatives are connected for every \(p<1\), while every
tridiagonal fixed-protocol representative at \(p=1\) is disconnected.
\end{corollary}

\begin{proof}
Choose \(n=2q+1\) above both the moment order and the exponent required by
the sine-power bound. Use the cyclic construction for \(p<1\) and the
block construction for \(p=1\). Both are unitary conjugates of the same
diagonal Hamiltonian and map the same fixed first site to the appropriate
spectral state. All amplitude, moment and window identities are unchanged.
\end{proof}

\section{Occupation and fermionic-sector lifts}

Let \(J\) be any Hermitian tridiagonal matrix on the one-particle site
space. On the occupation basis labelled by subsets \(S\subseteq
\{0,\ldots,n\}\), define a number term for each diagonal entry and allow a
directed hop \(j\to i\) only when \(i,j\) are adjacent, \(j\in S\), and
\(i\notin S\). Add the Hermitian transpose and divide by two.

\begin{theorem}[Exact local occupation embedding]
The resulting occupation matrix is Hermitian, has zero entries between
subsets of different cardinality, and its singleton block is exactly
\(J\):
\[
 X_{\{i\},\{j\}}=J_{ij}.
\]
\end{theorem}

\begin{proof}
Every number term is diagonal. Every hop erases one occupied site and
inserts one unoccupied adjacent site, so it preserves cardinality.
Hermitian symmetrization gives \(X^\dagger=X\). On singleton occupations,
the number terms return \(J_{ii}\), while the unique allowed hop from
\(\{j\}\) to \(\{i\}\) returns \(J_{ij}\); tridiagonality kills every
remaining off-diagonal entry.
\end{proof}

For the complete noninteracting spectrum, work intrinsically on
\(\bigwedge^k\C^{n+1}\). Define \(\mathrm d\Gamma(J)\) by applying \(J\)
in each tensor slot, summing the \(k\) results, antisymmetrizing, and
dividing by \(k!\).

\begin{theorem}[Complete spectrum in every fermion sector]
Suppose \(Ju_j=E_ju_j\) and \((u_j)_{j=0}^{n}\) is a basis. For every
\(k\)-element subset \(S=\{j_1<\cdots<j_k\}\),
\[
 \mathrm d\Gamma(J)(u_{j_1}\wedge\cdots\wedge u_{j_k})
 =
 \left(\sum_{j\in S}E_j\right)
 (u_{j_1}\wedge\cdots\wedge u_{j_k}).
\]
These wedges form a basis of \(\bigwedge^k\C^{n+1}\).
\end{theorem}

\begin{proof}
Before antisymmetrization, applying \(J\) in slot \(\ell\) multiplies the
wedge by \(E_{j_\ell}\). Summing the slots gives
\(\sum_{\ell}E_{j_\ell}\). Under each permutation, the exterior product
and the antisymmetrizer contribute the same sign; their product is one.
The \(k!\) equal terms cancel the normalization. Exterior products of
increasing \(k\)-subsets of a one-particle basis are the standard Slater
basis.
\end{proof}

\begin{corollary}[Common full Fock spectrum]
For every fixed-protocol candidate \(\widetilde K_\theta\), every
fermion-number sector has the complete eigenvalue list
\[
 \left\{\sum_{j\in S}\frac{\pi j}{T}:
 S\subseteq\{0,\ldots,n\},\ |S|=k\right\},
\]
with the same multiplicities for every \(\theta\in[0,1]\). For every
\(\theta\), its one-particle generator is a nearest-neighbour Hamiltonian
of Section 4; the endpoint generator is the active--spectator block form.
\end{corollary}

\begin{proof}
The one-particle eigenvectors are
\(U_\theta^\dagger|j\rangle\), with eigenvalues \(\pi j/T\). Apply the
preceding theorem. The subset index supplies the complete basis and counts
multiplicity.
\end{proof}

\section{Closed amplitude formulas}

Set \(z=\ee^{\ii\pi t/T}\) and define
\[
 A_0(t)=\sum_{j=0}^{n}q_jz^j,\qquad
 A_1(t)=\sum_{j=0}^{n}r_jz^j.
\]

\begin{lemma}[Two generating functions]
\begin{align}
 A_0(t)&=\left(\frac{1+z}{2}\right)^n,
\label{si:eq:A0}\\
 A_1(t)&=\left(\frac{1+z}{2}\right)^n+
          \left(\frac{z-1}{2}\right)^n.
\label{si:eq:A1}
\end{align}
\end{lemma}

\begin{proof}
The first identity is the binomial theorem. For the signed contribution to
the second,
\[
\sum_{j=0}^{n}q_j(-1)^{n-j}z^j
=2^{-n}(z-1)^n.
\]
Adding the unsigned contribution proves the result.
\end{proof}

The clock with target return probability \(p\in[0,1]\) is
\(\theta=\sqrt p\). Its amplitude is
\[
 A_p(t)=(1-\sqrt p)A_0(t)+\sqrt p\,A_1(t).
\]

\begin{theorem}[Every exact terminal probability]
For every \(p\in[0,1]\),
\begin{align}
 A_p(t)&=A_0(t)+\sqrt p
   \left(\frac{\ee^{\ii\pi t/T}-1}{2}\right)^n,
\label{si:eq:Ap}\\
 A_p(T)&=\sqrt p\,(-1)^n,\qquad |A_p(T)|^2=p.
\label{si:eq:target}
\end{align}
\end{theorem}

\begin{proof}
Equation~\eqref{si:eq:Ap} follows from the generating functions. At
\(t=T\), \(z=-1\): the first term is zero and the second is
\(\sqrt p(-1)^n\). Squaring its modulus gives \(p\).
\end{proof}

\section{Exact moment matching}

\begin{lemma}[Alternating binomial moments]
For every integer \(k\) with \(0\leq k<n\),
\[
\sum_{j=0}^{n}(-1)^{n-j}{n\choose j}j^k=0.
\]
\end{lemma}

\begin{proof}
The left side is the \(n\)th forward difference at zero of the polynomial
\(x\mapsto x^k\). An \(n\)th forward difference annihilates every
polynomial of degree below \(n\).
\end{proof}

\begin{theorem}[Common energy moments]
For every \(0\leq k<n\), every \(p\in[0,1]\), and
\(E_{j,T}=\pi j/T\),
\[
\sum_{j=0}^{n}w_{\sqrt p,j}E_{j,T}^{\,k}
=\sum_{j=0}^{n}q_jE_{j,T}^{\,k}.
\]
\end{theorem}

\begin{proof}
The difference between \(r_j\) and \(q_j\) is
\[
q_j(-1)^{n-j}=2^{-n}{n\choose j}(-1)^{n-j}.
\]
The preceding lemma annihilates its \(k\)th moment. Multiplication by the
common factor \((\pi/T)^k\) preserves equality. Convex interpolation between
\(q\) and \(r\) preserves it again.
\end{proof}

\section{Exponential finite-window agreement}

\begin{lemma}[Exact amplitude discrepancy]
For every \(p\in[0,1]\) and \(t\in\R\),
\[
 |A_p(t)-A_0(t)|
 =\sqrt p\left|\sin\left(\frac{\pi t}{2T}\right)\right|^n.
\]
\end{lemma}

\begin{proof}
Equation~\eqref{si:eq:Ap} and
\[
\left|\frac{\ee^{\ii x}-1}{2}\right|
=\left|\sin\frac{x}{2}\right|
\]
give the identity.
\end{proof}

\begin{lemma}[Probability is Lipschitz in amplitude]
If \(|a|\leq1\) and \(|b|\leq1\), then
\[
\bigl||a|^2-|b|^2\bigr|\leq2|a-b|.
\]
\end{lemma}

\begin{proof}
The reverse triangle inequality gives
\[
\bigl||a|-|b|\bigr|\leq|a-b|.
\]
Multiply by \(|a|+|b|\leq2\).
\end{proof}

Every \(A_p\) is a convex characteristic function of a probability measure,
so \(|A_p(t)|\leq1\). Let
\[
 P_p(t)=|A_p(t)|^2.
\]

\begin{theorem}[Uniform continuum bound]
If \(0\leq H<T\), then
\begin{align}
\sup_{\substack{p\in[0,1]\\|t|\leq H}}
 |A_p(t)-A_0(t)|
&\leq
\left[\sin\left(\frac{\pi H}{2T}\right)\right]^n,
\label{si:eq:ampwindow}\\
\sup_{\substack{p\in[0,1]\\|t|\leq H}}
 |P_p(t)-P_0(t)|
&\leq
2\left[\sin\left(\frac{\pi H}{2T}\right)\right]^n.
\label{si:eq:probwindow}
\end{align}
\end{theorem}

\begin{proof}
For \(|t|\leq H<T\), the absolute sine is monotone after reduction to the
interval \([0,\pi/2]\), and
\[
\left|\sin\left(\frac{\pi t}{2T}\right)\right|
\leq\sin\left(\frac{\pi H}{2T}\right)<1.
\]
Use \(\sqrt p\leq1\) for the amplitude bound, then apply the Lipschitz
lemma.
\end{proof}

\begin{corollary}[One record, every exact future]
For every \(T>0\), \(0\leq H<T\), \(\delta>0\), and \(m\in\mathbb N\),
there is a finite \(n>m\) such that, simultaneously for every \(p\in[0,1]\),
\[
\sup_{|t|\leq H}|P_p(t)-P_0(t)|<\delta,
\qquad P_p(T)=p,
\]
and the energy moments through order \(m\) equal those of the baseline
clock.
\end{corollary}

\begin{proof}
Put \(s=\sin(\pi H/(2T))\). Since \(0\leq s<1\), choose \(n>m\) with
\(2s^n<\delta\). Apply the common-moment theorem, the uniform continuum
bound, and equation~\eqref{si:eq:target}.
\end{proof}

\section{Minimum spectral support}

\begin{theorem}[Vandermonde moment separation]
Let \(x_0,\ldots,x_{r-1}\) be distinct real numbers and let
\((u_j)\), \((v_j)\) be two real weight lists. If
\[
\sum_{j=0}^{r-1}u_jx_j^k
=\sum_{j=0}^{r-1}v_jx_j^k
\quad\text{for }k=0,\ldots,r-1,
\]
then \(u_j=v_j\) for every \(j\).
\end{theorem}

\begin{proof}
The moment equations state \(V(u-v)=0\), where
\(V_{k,j}=x_j^k\). Its determinant is
\[
\det V=\prod_{0\leq i<j<r}(x_j-x_i)\neq0.
\]
Thus \(V\) is invertible and \(u-v=0\).
\end{proof}

\begin{corollary}[Sharp level count]
Two different spectral measures agreeing in their first \(n\) moments
(orders \(0,\ldots,n-1\)) require at least \(n+1\) distinct energy levels.
The pair \(q,r\) attains this lower bound.
\end{corollary}

\begin{proof}
At most \(n\) levels would be separated by the first \(n\) moments using the
theorem (or by its restriction to the actual number of levels). The
energies \(E_{j,T}=\pi j/T\), \(j=0,\ldots,n\), are distinct; \(q\neq r\);
and the exact moment-matching section proves equality of the required
moments.
\end{proof}

\section{Minimax and exact-data boundaries}

\begin{theorem}[One-half minimax lower bound]
Let \(F\) be any real-valued rule applied to the common record \(P_0\).
Under the conditions of the continuum corollary,
\[
\max\{|F(P_0)|,|F(P_0)-1|\}\geq\frac12.
\]
\end{theorem}

\begin{proof}
Write \(y=F(P_0)\). Since the compatible terminal set includes zero and
one,
\[
\max\{|y|,|y-1|\}\geq\frac12.
\]
\end{proof}

The construction uses dimension \(n+1\) and maximum angular energy
\(\pi n/T\). It does not claim ambiguity under simultaneous fixed bounds on
dimension and bandwidth. The result concerns every strictly positive error
radius; as that radius decreases or \(H/T\) approaches one, a larger finite
system is selected.

\end{document}
